% Generated draft: requires mathematical and source review.
\documentclass[11pt]{article}
\usepackage[margin=25mm]{geometry}
\usepackage{fontspec}
\setmainfont{Latin Modern Roman}
\usepackage{amsmath,amssymb,mathtools,unicode-math}
\setmathfont{Latin Modern Math}
\usepackage{enumitem,needspace,xurl,hyperref}
\setlength{\parindent}{0pt}
\setlength{\parskip}{5pt}
\setlength{\emergencystretch}{3em}
\newcommand{\N}{\mathbb{N}}
\newcommand{\Z}{\mathbb{Z}}
\newcommand{\Q}{\mathbb{Q}}
\newcommand{\R}{\mathbb{R}}
\newcommand{\C}{\mathbb{C}}
\newcommand{\F}{\mathbb{F}}
\newcommand{\A}{\mathbb{A}}
\newcommand{\PP}{\mathbb{P}}
\newcommand{\E}{\mathbb{E}}
\newcommand{\eps}{\varepsilon}
\newcommand{\dd}{\,\mathrm{d}}
\begin{document}
\section*{UnsolvedMath ID: 20002307}
\textbf{Source problem code:} catalogue entry
\subsection{Definitions and mathematical statement}
\textbf{Question 5 and missing context.} Must the operation $\oplus$ be definable in some $\Q\times F$, with $F$ finite? The source sentence does not define $\oplus$, the cover, allowed parameters or the structure on $\Q$. The aid recovers a strongly minimal group and a three-sheeted set map to the rational additive group from the preceding passage. The following formalization is conditional on that recovered setting; it is not an arbitrary assertion about every binary operation on a countable set.

Let $G$ be the intended strongly minimal group: every definable subset of its underlying set in elementary extensions is finite or cofinite. Let $\pi:G\to\Q$ be a three-to-one set map and choose a fiber labeling $h:G\to\Q\times F$, $|F|=3$, with first coordinate $\pi$. Transport $\oplus$ along $h$. Ask whether there exists a labeling/presentation in which the operation graph is definable in the pure additive structure $(\Q,+,0)$ together with finitely many named labels and permitted parameters. The pure-additive interpretation is an extra convention of the aid; a field or ordered structure on $\Q$ would give a different question.

\textbf{Supplied homomorphism obstruction.} The aid asserts that a surjective group homomorphism from a strongly minimal group to $(\Q,+)$ cannot have a nontrivial finite kernel. In particular a three-to-one map in the recovered context cannot be interpreted as such a homomorphism. This obstruction does not rule out a three-to-one \emph{set} map or prove that no definable finite-cover coordinate presentation exists.

\textbf{Supplied conditional coordinate reduction.} For each $i,j,k\in F$, define
\[
 \Gamma_{ij}^k=\{(x,y,z)\in\Q^3:(x,i)\oplus(y,j)=(z,k)\}.
\]
For the fixed labeling, definability of the transported operation is equivalent to definability of these finitely many coordinate graphs. If they are definable in the pure additive rational structure, its linear geometry gives one generic rule for each input-label pair,
\[
 (x,i)\oplus(y,j)
     =(a_{ij}x+b_{ij}y+c_{ij},\mu(i,j)),
 \qquad a_{ij},b_{ij}\in\Q\setminus\{0\}.
\]
Here generic means outside a lower-rank definable exceptional subset of $\Q^2$; parameters can supply the constants $c_{ij}$. Those exceptional loci and their operation values must be retained. There are nine generic rules for the three-sheeted case. Group cancellation forces the two slopes nonzero; these generic rules alone need not specify the full operation.

The aid's five generic associativity identities, for all $i,j,k\in F$, can be expressed as
\[
 \mu(\mu(i,j),k)=\mu(i,\mu(j,k)),
\]
\[
 a_{\mu(i,j),k}a_{ij}=a_{i,\mu(j,k)},\qquad
 a_{\mu(i,j),k}b_{ij}=b_{i,\mu(j,k)}a_{jk},
\]
\[
 b_{\mu(i,j),k}=b_{i,\mu(j,k)}b_{jk},
\]
\[
 a_{\mu(i,j),k}c_{ij}+c_{\mu(i,j),k}
       =b_{i,\mu(j,k)}c_{jk}+c_{i,\mu(j,k)}.
\]
They are necessary finite generic data for such a presentation, not sufficient to settle all exceptional cases or construct the requested labeling. No unconditional positive or negative answer to Question 5 is supplied.

\textbf{Separate Section 1.3, ``Continuous sections (Miller)''.} The appended paragraph is not a hypothesis of Question 5. It states that for a stable theory $T$ and $M\prec N\models T$, restriction $S_1(N)\to S_1(M)$ has a continuous section, in the Stone topologies, assigning each type its unique nonforking extension. Its connection to the preceding finite-cover question is not asserted by this extract.
\subsection{Short English statement}
The isolated question asks for finite-cover coordinates in which a strongly minimal group's operation becomes definable. The aid excludes a three-to-one homomorphism interpretation and gives necessary generic affine rules under a chosen pure-additive rational presentation. It leaves both the missing source conventions and the lower-rank exceptional behavior unresolved; the Miller type-space paragraph belongs to a separate section.
\subsection{Sources}
\begin{itemize}
\item Local UnsolvedMath registry, permanent ID 20002307, source code catalogue entry; canonical \texttt{aim-logic-notes.json}, notes[82].
\item AIM \emph{Computable stability theory}, Question 5 at the end of Section 1.2 and the appended Section 1.3 ``Continuous sections (Miller)'': \url{https://aimath.org/pastworkshops/computestabproblems.pdf}.
\item The supplied verification reports recovery from p.~2 of the original problem-session PDF. Its full preceding passage is not included here; the strongly minimal three-cover interpretation comes from the aid. No separate proving citation is supplied for its homomorphism or affine-cover claims.
\item Registry research attempt 1, dated 2026-08-14: partial result, status at run unknown, rigor ``proved'', novelty confidence low. The explicitly machine-generated, non-peer-reviewed aid requires independent verification before citation.
\end{itemize}
\end{document}
